\section{The canonical conformal structure}
\label{conf:section}

We identify the conformal structure obtained by gluing the fixed flag
triangles.  There are two steps.  First, extremal length cannot degenerate:
a degenerating annulus would produce a nonzero flow of rectifiable curves
in the reference coordinate, and locality would then produce cheap
crossings in two transverse directions.  Second, the resulting
quasiconformal comparison has a local infinitesimal conformal structure.
The same locality argument makes this structure deterministic and round.
Neither step uses the intrinsic metric comparison.

Throughout this section, a reference coordinate means the original
oriented peanosphere coordinate of Section~\ref{enc:section}.  On the
finite sphere let
\(H_n:X_n\to\Chat\) be the homeomorphisms of
Theorem~\ref{top:homeomorphism}, chosen from the contour and surface
data before drawing the three fresh normalization samples.  On a compact
planar test domain we use their protected-window counterparts from
Proposition~\ref{top:protected-window}.  In particular,
\begin{equation}
 \max_T\diam H_n(T)\longrightarrow0
 \label{conf:topological-mesh}
\end{equation}
in the reference coordinate.  This assertion is topological and precedes
the uniformization of \(X_n\).

All extractions below retain the compatible local arrays of
Theorem~\ref{ker:local-kernels}: actual flag-surface extremal lengths,
strict port and confinement tests, and their joint refinements.  A test
whose projected confinement is compactly inside an open set can be
specified by a finite union of protected cells with an intermediate
margin.  We always fix that finite test before passing to the discrete
limit.  The kernel theorem permits countably many such tests and
arbitrary deterministic positive units for extremal length.  Local
changes of ordinary field law use Proposition~\ref{ker:field-transfer}.
Transfer to the finite sphere at its actual area-normalized heights uses
Lemma~\ref{ker:sphere-density} together with the common kernels. For
probability-one conclusions its domination direction suffices.

\subsection{Extremal length and flows}
\label{conf:el-preliminaries}

We use the extremal-length and electrical-energy framework of
Ahlfors and Beurling \cite{AhlforsBeurling1952ConformalInvariants}.
For a family \(\mathcal C\) of curves on a Riemann surface, we use
\begin{equation}
 \EL(\mathcal C)=
 \sup_{\sigma}
 \frac{\bigl(\inf_{\gamma\in\mathcal C}
                 \int_\gamma\sigma\,\dd s\bigr)^2}
      {\int\sigma^2\,\dd A}.
 \label{conf:el-definition}
\end{equation}
Here \(\dd s,\dd A\) are any conformal background length and area
elements and the supremum is over nonnegative measurable densities of
finite positive energy.  Equivalently, \(\EL\) is the reciprocal of
the infimal energy of a density charging every member of \(\mathcal C\)
by at least one.  On \(X_n\) the background representative is the
actual Euclidean flag geometry.  The extended charts at the finitely
many conical vertices give the same definition.

For an annulus \(A\), write \(\EL_{\rm tr}(A)\) for its through
crossings and \(\EL_{\rm wind}(A)\) for its nonzero-winding closed
curves.  Uniformization by a cylinder, and by a rectangle for a
quadrilateral, gives
\begin{equation}
 \EL_{\rm tr}(A)\EL_{\rm wind}(A)=1,
 \qquad
 \EL(\mathcal C_{\leftrightarrow}(R))
 \EL(\mathcal C_{\updownarrow}(R))=1.
 \label{conf:duality}
\end{equation}
For nested disjoint annuli crossed in series, through extremal lengths
add as a lower bound.  One can see this directly by adding admissible
densities with disjoint supports and optimizing their relative weights.

A \emph{unit flow} is a probability measure \(\nu\) on rectifiable
curves.  Its traffic is the measure obtained by averaging arclength,
with multiplicity.  If this measure has density \(\tau\) with respect
to \(\dd A\), then
\begin{equation}
 \int\!\left(\int_\gamma\sigma\,\dd s\right)\nu(\dd\gamma)
   =\int\sigma\tau\,\dd A,
 \qquad
 \EL(\mathcal C)\le\int\tau^2\,\dd A
 \quad\hbox{if }\nu(\mathcal C)=1.
 \label{conf:flow-bound}
\end{equation}
Indeed, the first identity is Tonelli's theorem.  For a density
admissible for \(\mathcal C\), the left-hand side is at least one;
Cauchy--Schwarz and the reciprocal variational formula prove the second
assertion.  Restricting the curves to subpaths decreases traffic.
Restricting to an event of flow probability \(u>0\) and normalizing
multiplies the traffic energy by at most \(u^{-2}\).

For a cylinder \([0,L]\times(\R/H\Z)\), choose a horizontal segment
with its second coordinate uniform.  Its traffic density is \(H^{-1}\),
and its energy is \(L/H=\EL_{\rm tr}\).  Thus every conformal annulus
has a unit through flow of energy equal to its through extremal length.
Trimming between two intermediate separating curves gives compactly
confined subpaths with no increase of energy.  These observations also
apply to piecewise smooth annuli in a flag surface.

\subsection{Crossing quantiles and the two local tests}
\label{conf:quantile-subsection}

Fix a disk compactly inside an ordinary part of the cone coordinate,
away from its marked singularity.  Inside this disk choose a center
\(z_0\) and radii
\[
 0<r_0<r_1<r_2<R_2<R_1<R_0,
\]
with additional strict collars around the six circles.  These give a
wide traversal band, a middle band \(A(z_0,r_1,R_1)\), and a narrower
band between radii \(r_2,R_2\).  All radii and collars remain fixed.

Use the bilateral word with microscopic time unit \(4^j\), hence
diffusive height unit \(2^j\).  The protected-window construction of
Proposition~\ref{top:protected-window} gives a finite flag surface
containing these tests.  In it choose a piecewise smooth annulus
\(\mathcal A_j\) approximating the inverse image of the middle
band under the topological correspondence, with its two boundary curves
separating the corresponding inner and outer collars.  Make this a
fixed measurable choice: enumerate annuli with rational polygonal
boundaries in finite subdivisions of the flag complex, and take the
first satisfying these separation tests with boundary-image error
less than \(2^{-j}\) from the two middle-band circles.  Polygonal
approximation of their inverse Jordan curves gives such annuli;
measurability of the correspondence makes the tests measurable.  Put
\(Z_j=\EL_{\rm tr}(\mathcal A_j)\).  If this temporary construction
fails, put \(Z_j=1\); its failure probability tends to zero.

The chosen annulus need not itself be a local readout.  Its two uses
are captured by local curve families with strict ports and confinement.
Let \(\mathcal W_j\) be the wider crossing family, from the inner
port near radius \(r_0\) to the outer port near radius \(R_0\),
confined to a strict annular collar of \(A(z_0,r_0,R_0)\).
Every member contains a through crossing of \(\mathcal A_j\).
A path from the central hole to the exterior of the outer buffer
has such a wider crossing as a subpath: stop on its first outer
visit and start at its last preceding inner visit.  Conversely, trim each curve
of the unit cylinder flow on \(\mathcal A_j\) between its last visit
to radius \(r_2\) before its first subsequent visit to radius \(R_2\).
The trimmed pieces lie in the closed band between those circles.
They belong to a narrower local family \(\mathcal N_j\), whose
ports and confinement include strict collars of that band.  Their
projected endpoints have a fixed positive separation.  If \(\tau_j\)
is its traffic density with respect to Euclidean flag area
\(\dd A_j\), then
\begin{equation}
 \EL(\mathcal W_j)\ge Z_j,
 \qquad
 \EL(\mathcal N_j)\le\int\tau_j^2\,\dd A_j\le Z_j.
 \label{conf:band-sandwich}
\end{equation}
For the first inequality, extend a density on \(\mathcal A_j\) by
zero; every wider crossing pays for a through subpath.  The second
uses \eqref{conf:flow-bound} and the decrease of traffic under trimming.

Implement the two families by intermediate protected-cell covers of
the ports and confinements.  The extra collars allow the wider ports
to remain on opposite sides of \(\mathcal A_j\), and the narrower
cover to contain every trimmed piece.  The full-preimage protection
and local readouts of Lemma~\ref{ker:readouts} supply such finite
covers, with margins fixed before taking the discrete limit.  Their
countable refinements are retained together.  Thus the inequalities
in \eqref{conf:band-sandwich}, on the event that the construction
succeeds, transfer its quantile bounds to local EL tests.  They require
no locality of the chosen topological homeomorphism or annulus.

Fix a sufficiently small \(\varepsilon>0\), and choose positive
lower quantiles \(\beta_j\) satisfying
\begin{equation}
 \PP[Z_j\le2\beta_j]\ge\varepsilon/2,
 \qquad
 \PP[Z_j\ge\beta_j/2]\ge1-2\varepsilon.
 \label{conf:quantiles}
\end{equation}
These inequalities allow atoms.  The second gives high-probability
barriers for the wider crossing family; the first gives a
positive-probability unit flow in the narrower family with energy
at most \(2\beta_j\).  The vanishing construction failures are
absorbed in the probability slack.  Choose \(\varepsilon\) so small
that the lower test, after the fixed local changes of field law and
band margins, meets the probability threshold of
Theorem~\ref{ann:aligned-annuli}.  The failure probabilities under
these fixed changes of density tend to zero with \(\varepsilon\).

We will rule out \(\liminf_j\beta_j=0\).  On a hypothetical
record-low sequence \(j_m\), set \(\beta_m=\beta_{j_m}\).
For every fixed backward offset \(i\ge0\),
\(\beta_{j_m-i}\ge\beta_m\) eventually, so the wider-band tests
still give barriers in \(\beta_m\) units.  Their time units
\(4^{j_m-i}\) have relative labels \(v_i=4^{-i}\), exactly the
labels used in annular alignment.  Extremal length has no additional
power of \(v_i\).  The next two lemmas turn these barriers and the
small-energy flow at offset zero into a nonzero absolutely continuous
average arclength measure in the reference coordinate.  Local germs
will then force cheap crossings in two transverse directions.

\subsection{A density that charges all projected advances}
\label{conf:density-subsection}

We record precisely the consequence of the aligned-annulus construction
that is needed here.  Suppose that in an extraction, indexed by \(m\),
the high-probability shell tests give a lower through extremal length
of order \(\beta_m>0\), at the backward labels required by
Theorem~\ref{ann:aligned-annuli}.  The units \(\beta_m\) are
deterministic; the record-low experiment above is the first application.
The aligned-annulus result of \Cref{ann:aligned-annuli}, applied to
these local tests, supplies the following configuration for arbitrarily
large integers \(N\), with summable failure probabilities after
subextraction.  There are at most \(C N\) distinct geometric radii
between \(e^{-BN}\) and \(e^{-aN}\).  At radius \(r\) the centers
form a mesh of spacing comparable to \(r\), and their buffered shells
have bounded overlap.  Every point of a fixed compact set lies in the
central holes of at least \(cN\) good shells.  For a good shell, every
path from its central hole to the exterior of its outer buffer is
charged by a density supported in that buffer, with charge at least
\(r\) and energy at most \(Cr^2/\beta_m\).  The constants depend
only on the fixed shell geometry and the probability cutoff.  All these
statements concern actual discrete-surface paths and supports under
the homeomorphism \(H_m\), with strict margins.

The density assertion follows from the extremal-length lower bound and
\eqref{conf:el-definition}, with a fixed slack in the constant.  If a
measurable choice is required, take a near-minimizing condenser
potential and its gradient.  Continuous piecewise affine potentials
on successive finite refinements of the flag triangles, with rational
values and fixed boundary values on the two port sets, form a
countable energy-dense family.  Selecting the first potential below
the prescribed energy threshold gives a measurable choice.  Its
gradient, enlarged on triangle edges by the larger one-sided value,
charges every crossing; this change on edges does not change energy.

\begin{lemma}[Simultaneous density budget]
\label{conf:density-budget}
Let \(U\) be a finite union of rational boxes compactly inside an
ordinary reference chart.  In the preceding configuration put
\(R_N=e^{-aN}\).  There are nonnegative densities \(g_{m,N,U}\),
supported over \(U^{+CR_N}\), such that
\begin{align}
 \limsup_{m\to\infty}\beta_m
       \int g_{m,N,U}^{2}\,\dd A_m
       &\le C\Area(U^{+CR_N}),
       \label{conf:density-energy}\\
 \int_{\gamma|_I}g_{m,N,U}\,\dd s_m
       &\ge c\diam(H_m\gamma(I))-CR_N-o_m(1)
       \label{conf:density-charge}
\end{align}
for every path interval \(I\) projected into \(U\).
For any fixed finite number of pairwise disjoint path intervals the
charges add.  The same density works for every path and every such
collection.  The constants are independent of \(U,N,m\); the
threshold in \(m\) may depend on the finite configuration and margins.
\end{lemma}

\begin{proof}
At each center and radius retain one good shell density.  Sum them over
the centers whose buffers meet \(U\), and then sum over the retained
radii and divide by \(N\).  At radius \(r\), bounded overlap and
grid counting give total energy at most
\(C\Area(U^{+CR_N})/\beta_m\).  There are at most \(CN\) radii.
Cauchy--Schwarz for their sum proves
\eqref{conf:density-energy}.

For the charge, choose a coordinate in which the path interval has
oscillation at least its diameter divided by \(\sqrt2\), and restrict
to a subinterval between a minimum and a maximum of that coordinate.
Discard coordinate values within \(CR_N\) of either endpoint.  At
each remaining value choose a visit of the path.  This visit belongs
to at least \(cN\) good central holes.  A hole at radius \(r\) can
account for a coordinate interval of length at most \(Cr\), and the
path must make a full crossing of its buffered shell.  Integrating the
number of covering holes over the coordinate interval gives
\[
 cN\bigl(\operatorname{osc}-CR_N\bigr)
       \le C\sum_{\text{shells crossed}}r.
\]
The integral of the summed shell densities is at least the sum on the
right before the harmless fixed factor.  Divide by \(N\).  Strict
inner and outer margins absorb the uniform projection errors, proving
\eqref{conf:density-charge}.  Integrals over disjoint parameter
intervals add even when their traces overlap.  No choice of the density
was made using the path.
\end{proof}

\subsection{Compact flows and average arclength}
\label{conf:flows}

The density budget now controls how many disjoint advances a flow can
make.  This yields compactness without assuming that the projected
discrete curves have bounded Euclidean length.

\begin{lemma}[Flow extraction]
\label{conf:flow-compactness}
Suppose that the configurations of Lemma~\ref{conf:density-budget}
hold along a sequence for all sufficiently large members of a
deterministic sequence of \(N\)'s, simultaneously for the countable
box library.  Let \(\nu_m\) be unit flows of compactly confined paths
whose projected endpoints are separated by at least \(d_0>0\), and
whose traffic satisfies
\[
 \int\tau_m^2\,\dd A_m\le C_0\beta_m.
\]
After reparameterization and subextraction, the projected flow laws
converge weakly to a probability law \(\nu\) on continuous paths.
Its paths have endpoint separation at least \(d_0\), are rectifiable
almost surely, and their average arclength measure \(\Lambda\)
satisfies
\begin{equation}
 \Lambda(U)\le C\sqrt{\Area(U)}
 \label{conf:average-length}
\end{equation}
for every finite union of rational boxes compactly in the chart.
In particular, \(0<\Lambda(K)<\infty\) on a sufficiently large compact
confinement \(K\), and \(\Lambda\) is absolutely continuous with
respect to planar area.
\end{lemma}

\begin{proof}
Fix \(h>0\), and choose \(N\) with \(CR_N<ch/4\).  A collection
of disjoint path intervals each of projected diameter at least \(h\)
then pays at least \(ch/2\) per interval for all sufficiently large
\(m\).  By \eqref{conf:flow-bound} and
\eqref{conf:density-energy}, its expected cardinality is bounded by
\(C/h\).  The same bound holds for the maximal cardinality, by
monotone convergence over finite collections.

Here is a measurable reparameterization implementing compactness.
For a continuous projected path, mark its successive first moves by
distance \(2^{-k}\), and denote their finite number by \(J_k\).
Give each such mark clock mass \(2^{-k}/(1+J_k)\), for every
\(k\ge1\), and add normalized Lebesgue mass in its original
parameter.  The total clock mass lies in \([1,2]\).  Use its inverse,
inserting a pause at each atom, and normalize the new parameter to
\([0,1]\).  This is a representative of the same path up to zero
uniform reparameterization distance.  On a set where \(J_k\le M_k\)
for every \(k\), a new parameter interval shorter than
\(2^{-k-1}/(1+M_k)\) contains no complete move at scale \(2^{-k}\),
and the image oscillation there is at most \(2^{2-k}\).

The expected crossing-count bounds allow \(M_k<\infty\) to be
chosen with total exceptional probability as small as desired.  For
each fixed \(k\), the finitely many indices preceding its eventual
bound can be included by increasing \(M_k\): their crossing counts
are finite almost surely.  Thus the reparameterized laws are tight
in \(C([0,1],K)\), by Arzel\`a--Ascoli.  Prokhorov's theorem gives
\(\nu\), and endpoint separation passes to the limit.

To prove the length assertion without exchanging infinite suprema,
fix \(J\).  Consider at most \(J\) disjoint closed parameter intervals
with rational endpoints whose images have positive clearance inside
\(U\), and sum the distances between their endpoint images.  The
supremum of these sums is a nonnegative lower semicontinuous path
functional: each strict confinement test is open, and each finite sum
is continuous on that test.  For every such collection,
\eqref{conf:density-charge} and Cauchy--Schwarz bound its flow
expectation by
\[
 C\sqrt{\Area(U^{+CR_N})}+CJR_N+o_m(1).
\]
The bound also holds for the supremum, because the density charge
holds simultaneously for all collections.  Pass first to the weak
flow limit, then let \(N\to\infty\) with \(J\) fixed, and finally
let \(J\to\infty\).  These finite interval sums exhaust variation
in the open set \(U\).  Taking one larger confinement first proves
finite total variation for almost every path; the same argument in
each \(U\) then proves \eqref{conf:average-length} for its arclength
measure, with multiplicity.

Every limit path has length at least \(d_0\), so \(\Lambda\ne0\).
The constant in \eqref{conf:average-length} is uniform over finite
unions of boxes.  Cover a planar null set by countably many rational
boxes of arbitrarily small total area, apply the estimate to finite
subunions, and increase to the full union.  This proves
\(\Lambda\ll\dd z\).
\end{proof}

The lemma applies pathwise on a positive-probability event in a joint
extraction.  One may retain the flow laws as additional random
variables and use the preceding tightness argument after restricting
to compact bounds.  The local kernel theorem is not conditioned on
these flow laws: the flows will only witness events in the already
retained local extremal-length arrays.

\subsection{Local germs}
\label{conf:germ-subsection}

We next isolate the pointwise use of locality.  Write
\(\mathcal H\) for the sigma field of the real and imaginary fields
in an ordinary chart.  A \emph{local germ readout} \(Y(z)\) is a
jointly measurable function of the retained local arrays with the
following property: for every neighborhood \(V\) of \(z\), it is
readable from strict inner and outer tests in arbitrarily small
subneighborhoods of \(z\) contained in \(V\).  Countable finite
port covers are used in this definition.

At almost every deterministic \(z\), Lemma~\ref{ker:single-traversal}
says that all visits to a sufficiently small neighborhood of \(z\)
belong to one traversal in any prescribed larger neighborhood.
The needed success event is local: fix two surrounding buffers inside
that larger neighborhood. By Lemma~\ref{enc:local-exploration}, the
imaginary field there determines the entire finite list of directed
outer-buffer traversals meeting a prescribed compact inner disk. The
condition that this list has cardinality one is therefore measurable
there. Shrinking the inner disk exhausts a probability-one event by
unique visitation and properness, as in Lemma~\ref{ker:single-traversal}.
Consequently the germ requires no exterior ordering choice. The
kernel theorem and its restriction compatibility imply that
\(\EE[\varphi(Y(z))\mid\mathcal H]\) is measurable with respect
to the restrictions of the fields to every such larger neighborhood.
For two distinct ordinary points these readouts have independent
conditional kernels.  These conclusions include local auxiliary marks,
which may be integrated out.  The single-traversal exhaustion has
probability one; its success in an inner neighborhood is testable in
the larger neighborhood, as stipulated in
Lemma~\ref{ker:single-traversal}.

\begin{lemma}[Deterministic germs and rotations]
\label{conf:germs}
A bounded local germ readout with values in a separable metric space
is, for area-almost every deterministic point, almost surely
deterministic.  Its deterministic value respects every coordinate
rotation under which the intrinsic readout transforms.  The statements
hold jointly for countably many readouts and rotations.
\end{lemma}

\begin{proof}
First take a bounded real readout and set
\(p(z)=\EE[Y(z)\mid\mathcal H]\), with jointly measurable versions.
Conditional independence at distinct points and Fubini give, for
every rational disk \(V\),
\[
 \EE\left[\left(\int_V(Y(z)-p(z))\,\dd z\right)^2
              \middle|\mathcal H\right]=0.
\]
The diagonal in \(V\times V\) has area zero.  A countable disk
exhaustion and Lebesgue differentiation imply
\(Y(z)=p(z)\) for almost every \(z\), almost surely.  A second
application of Fubini gives this identity almost surely at almost
every deterministic point.

For completeness, the point germ of a Dirichlet GFF is trivial.
Realize the field as an isonormal Gaussian process over
\(H^1_0(V)\), with its Dirichlet inner product.  The Gaussian
subspace generated by its restriction to \(B(z,r)\) is the closed
span of the Green potentials of test functions supported there.
An element in the intersection of these subspaces is orthogonal to
every smooth function vanishing near \(z\).  Such functions are
dense in \(H^1_0(V)\): a logarithmic cutoff around \(z\) has
Dirichlet energy tending to zero, so a point has zero capacity.
The subspace intersection is therefore zero.  Decreasing orthogonal
projections, first on each finite Gaussian chaos and then by
\(L^2\) approximation, show that the intersection of the completed
Gaussian sigma fields is trivial.  The same argument applies to the
direct sum of two independent fields.  Taking the imaginary field
modulo its period only decreases the observable sigma fields.
Local absolute continuity in Proposition~\ref{ker:field-transfer}
transfers triviality to the actual field restrictions and their phase
conventions.

For almost every deterministic \(z\), the conditional expectation
\(p(z)\) is measurable with respect to each shrinking field
restriction, hence to this trivial germ.  It is deterministic, and
so is \(Y(z)\).  A countable family of bounded continuous functions
separating points proves the assertion for a separable target space.

Finally rotate a local coordinate about \(z\), making the corresponding
imaginary-field phase adjustment.  The intrinsic time/height ports
and flag-surface variational outputs are unchanged; only their
spatial coordinate descriptions transform.  This is the affine
covariance in Proposition~\ref{ker:field-transfer}.  The original
and transformed field laws are locally equivalent.  Thus every
probability-zero or probability-one assertion about a germ transports
to the transformed germ.  Since its value is deterministic, that
value must obey the claimed rotation rule.  Taking countable
intersections completes the proof.
\end{proof}

\subsection{Extremal length does not degenerate}
\label{conf:nondegeneration-subsection}

\begin{proposition}[No degeneration of crossing quantiles]
\label{conf:no-degeneration}
There exists \(\beta_*>0\) such that
\(\liminf_{j\to\infty}\beta_j\ge\beta_*\).
\end{proposition}

\begin{proof}
Suppose otherwise, and take the record-low extraction described in
Section~\ref{conf:quantile-subsection}: successive record lows \(j_m\) with
\(\beta_{j_m}\to0\).  For every fixed backward offset \(i\),
\(\beta_{j_m-i}\ge\beta_{j_m}\) eventually.  Retain all fixed
offset laws and all local EL arrays in units
\(\beta_m=\beta_{j_m}\), using the diagonal extraction in
Theorem~\ref{ker:local-kernels}.  The second inequality of
\eqref{conf:quantiles}, through \eqref{conf:band-sandwich}, gives
the wider-band shell barriers in these units at every fixed backward
label \(v_i=4^{-i}\).  The aligned-annulus result and
Lemma~\ref{conf:density-budget} therefore apply.

The first inequality of \eqref{conf:quantiles} retains a
positive-probability small-EL event.  Retain the trimmed cylinder flows
from \eqref{conf:band-sandwich} as witnesses in the joint extraction.
On a fixed buffer containing their confinement and the local tests,
Proposition~\ref{ker:field-transfer} gives a strictly positive
Radon--Nikodym derivative of the ordinary field-pair restrictions
relative to the cone restrictions.  Reweight the entire coupled
extraction by this field-dependent derivative.  The local field and
array marginal then has the ordinary law with the same kernels;
the positive-probability flow event and its pathwise inequalities
persist under this equivalent change of law.  No kernel is conditioned
on the flows.  Their paths have fixed positive advance in a compact
set and traffic energy \(O(\beta_m)\).  Lemma~\ref{conf:flow-compactness}
gives, on that event, a probability flow \(\nu\) of nonconstant
rectifiable paths with nonzero absolutely continuous average arclength
\(\Lambda\).

We describe the local event witnessed by these paths.  Take a finite
sufficiently fine set of unoriented directions, closed under a
quarter-turn.  For a direction \(v\), a center \(z\), and
\(s>0\), put
\[
 R_v(z,s)=\{z+t v+u v^\perp: |t|<3s/2,\ |u|<s/4\}.
\]
A strict crossing runs from \(t<-s\) to \(t>s\), has compact
confinement in this rectangle, and is represented by intermediate
port covers.  Let \(E_v(z)\) be the event that at every sufficiently
small dyadic \(s\) some such cover has finite limiting EL in units
\(\beta_m\).  Explicitly, this is a countable union over the first
dyadic scale, followed by a countable intersection over smaller
scales, and a union over admissible finite covers and finite integer
EL bounds.  Whether a cover fits strictly inside the indicated
rectangle is a measurable open geometric test.  Thus \(E_v(z)\)
is jointly measurable in \(z\) and the local arrays.  Discarding
finitely many scales leaves the event unchanged, so it is a local
germ readout.  Its finite EL bound may depend on the scale.

At arclength-almost every interior point of almost every path of
\(\nu\), arclength parameterization has a unit tangent.  Choose
\(v\) within the angular slack of this tangent.  Differentiability
at that parameter value gives strict crossings of \(R_v(z,s)\)
at every sufficiently small \(s\).  Almost every path belongs to
the support of \(\nu\).  Each fixed strict tube crossing is an
open path event containing that path, so has positive \(\nu\)-mass.
Weak convergence gives positive mass in the discrete flows for all
sufficiently large \(m\).  Restrict and normalize that mass in
\eqref{conf:flow-bound}.  Its energy is still \(O(\beta_m)\), with
a finite constant depending on this radius and tube.  A finite
intermediate cover of positive mass supplies the required local EL
witness.  This proves \(E_v(z)\) at \(\Lambda\)-almost every
such point for some direction \(v\).

Since \(\Lambda\ne0\) and \(\Lambda\ll\dd z\), on the
positive-probability flow event at least one of the finitely many
events \(E_v\) holds on a set of positive planar area.
Lemma~\ref{conf:germs} makes its indicator deterministic at almost
every deterministic point and transports it under a quarter-turn.
By Fubini, on a probability-one event these equalities hold outside
one planar null set, simultaneously for the finite set of directions.
The positive-area flow witness therefore meets that full-area set,
where \(E_v(z)\) and \(E_{v^\perp}(z)\) have the same truth
value.  Choose a point satisfying both, and then one fixed dyadic
radius below their two starting scales.

Let \(\mathcal F_{v,m}\) and \(\mathcal F_{v^\perp,m}\) be the
two discrete curve families specified by the selected covers.
Their extremal lengths are at most \(C_1\beta_m\) and
\(C_2\beta_m\), respectively, for all sufficiently large \(m\).
Choose a common central quadrilateral \(Q_m\), with projected
sides near those of \(\{|t|<s/2,|u|<s/2\}\), and smooth it on
the discrete surface within the available margins.  As shown in
Figure~\ref{conf:transverse-figure}, the narrower strips have lateral
clearance from the inappropriate sides of this quadrilateral.  Every
member of either family therefore contains a subpath crossing the
corresponding pair of sides of \(Q_m\).  For any density \(\sigma\) on
\(Q_m\), its extension by zero to the surface has the same energy,
and
\[
 \inf_{\gamma\in\mathcal F_{v,m}}\int_\gamma\sigma\,\dd s_m
 \ge
 \inf_{\gamma\in\mathcal C_{\leftrightarrow}(Q_m)}
                       \int_\gamma\sigma\,\dd s_m.
\]
Taking the supremum in \eqref{conf:el-definition}, and arguing
identically in the perpendicular direction, yields
\[
 \EL(\mathcal C_{\leftrightarrow}(Q_m))
       \le\EL(\mathcal F_{v,m})\le C_1\beta_m,
 \qquad
 \EL(\mathcal C_{\updownarrow}(Q_m))
       \le\EL(\mathcal F_{v^\perp,m})\le C_2\beta_m.
\]
The constants are finite because the radius and the two covers
have now been fixed.  Their product
tends to zero, contradicting \eqref{conf:duality}.  This excludes the
record-low extraction.
\end{proof}

\begin{figure}[tb]
\centering
\begin{tikzpicture}[scale=.78,>=Latex]
  \fill[blue!7] (-3,-.5) rectangle (3,.5);
  \fill[orange!10] (-.5,-3) rectangle (.5,3);
  \draw[blue!55,dashed] (-3,-.5) rectangle (3,.5);
  \draw[orange!75!black,dashed] (-.5,-3) rectangle (.5,3);
  \draw[very thick] (-1,-1) rectangle (1,1);
  \node[anchor=south west] at (1,1) {$Q_m$};
  \draw[blue!75!black,thick,->]
       (-2.6,.08) .. controls (-1.4,-.3) and (-.6,.26) .. (0,.05)
                 .. controls (.9,-.22) and (1.5,.24) .. (2.6,-.05);
  \draw[orange!80!black,thick,->]
       (.12,-2.6) .. controls (-.28,-1.5) and (.24,-.7) .. (.04,0)
                 .. controls (-.2,.8) and (.24,1.5) .. (-.06,2.6);
  \node[anchor=west,text=blue!75!black] at (3.2,0)
       {$\EL_{\leftrightarrow}(Q_m)\le C_1\beta_m$};
  \node[anchor=south,text=orange!80!black] at (0,3.15)
       {$\EL_{\updownarrow}(Q_m)\le C_2\beta_m$};
  \node[anchor=north] at (0,-3.35)
       {$\EL_{\leftrightarrow}(Q_m)\,\EL_{\updownarrow}(Q_m)=1$};
\end{tikzpicture}

\caption{The dashed strips represent \(R_v(z,s)\) and
\(R_{v^\perp}(z,s)\).  Their strict lateral clearance forces
crossings of the corresponding sides of the common quadrilateral
\(Q_m\).  The picture uses the reference coordinate; the extremal lengths are those of the actual flag surface.
Finite bounds at one common radius contradict duality when
\(\beta_m\to0\).}
\label{conf:transverse-figure}
\end{figure}

\begin{proposition}[Macroscopic extremal-length bounds]
\label{conf:macroscopic-el}
In every diverging-scale extraction, including the finite sphere,
the following hold in the reference coordinate.
\begin{enumerate}
\item For every \(d>0\), the family of all discrete-surface paths
whose projected diameter is at least \(d\) has a positive limiting
lower extremal-length bound.
\item There is a deterministic \(c_*>0\) such that every fixed
round annulus \(A(z,r,2r)\) in an ordinary coordinate chart has
\[
 \liminf_n\EL_{X_n}\bigl(H_n^{-1}
                   \mathcal C_{\rm tr}(A(z,r,2r))\bigr)\ge c_*.
\]
The assertion holds simultaneously for a countable dense annulus
library, with strict margins.  It includes annuli surrounding an
isolated exceptional mark whose closures avoid the marks.
\end{enumerate}
\end{proposition}

\begin{proof}
Proposition~\ref{conf:no-degeneration} makes the high-probability
barriers uniformly positive in raw EL units. To pass from powers of four
to an arbitrary diverging integer sequence, write
$n=4^{j_n}v_n$, where $j_n=\lfloor\log_4 n\rfloor$ and $1\le v_n<4$,
and extract $v_n\to v\in[1,4]$ jointly with every fixed offset array.
Reexpress the same raw cuts and incidence readouts in the new units:
the area clock is divided by $v_n$, frontier heights by $\sqrt{v_n}$,
and the real field is shifted by $-\gamma^{-1}\log v_n$.
The actual flag-surface extremal length is unchanged. Strict chart
margins allow the converging factors to be retained through this
indexed-array identity. On each fixed buffered domain, the constant
shifts lie in a bounded family and admit cutoff extensions of uniformly
bounded Dirichlet norm. The two-sided Cameron--Martin comparison
therefore chooses a common high-probability barrier threshold before
annular alignment. This argument does not identify differently indexed
kernels at an unshifted field. Apply the aligned-annulus result and
Lemma~\ref{conf:density-budget} with \(\beta_m=1\).

For a curve making advance \(d\) compactly in \(U\), choose \(N\)
with \(CR_N<cd/4\).  The density charges the entire family by at
least \(cd/2\), for large \(m\), and has bounded energy.  Its
extremal length is therefore bounded below by a constant times
\(d^2/\Area(U^{+CR_N})\).  A finite collection of ordinary patches
catches a fixed positive subadvance of every curve of spherical
diameter at least \(d\).  To obtain such a collection, first remove
disjoint sufficiently small neighborhoods of the finitely many
exceptional marks; a curve of diameter \(d\) either makes an
advance away from them or crosses one of their surrounding compact
annuli.  Cover the resulting compact regions by finitely many
ordinary boxes and sum the densities.  This proves the first claim
for the entire family at once.

For a through crossing of \(A(z,r,2r)\), take a slightly larger
annular union of boxes of area at most \(Cr^2\).  Each crossing
makes projected advance at least a fixed multiple of \(r\) there.
The preceding quotient is bounded below by a universal constant.
The union can avoid the center and all other exceptional marks; in
particular this proves the assertion for fixed shells surrounding
an exceptional center.  Changes of ordinary field law preserve the
eventual statement and its constant.  Countable extraction handles
all rational tests and margins simultaneously. Lemma~\ref{ker:sphere-density}
transfers these probability-one conclusions to the finite sphere at its
actual area-normalized heights; Theorem~\ref{ker:local-kernels} retains
the same incidence and extremal-length readouts.
\end{proof}

\subsection{An all-center ring criterion}
\label{conf:analytic-subsection}

We give the analytic compactness statement in a form requiring no
estimate at a radius depending on \(n\).  In a finite conformal
atlas of the round sphere, fix the countable library of annuli
\(A(z,r,2r)\) with rational centers and radii and compact closure
in a chart.  Annuli around finitely many omitted centers are also
allowed, as in Proposition~\ref{conf:macroscopic-el}.

\begin{lemma}[Ring criterion]
\label{conf:ring-criterion}
Let \(f_n:\Chat\to\Chat\) be homeomorphisms.  Suppose that
\(f_n(z_n^i)=w_n^i\), where both triples converge to triples of
distinct points.  Suppose that, for every annulus \(A\) in the
preceding library,
\begin{equation}
 \liminf_n\EL_{\rm tr}(f_n(A))\ge c_*>0.
 \label{conf:ring-hypothesis}
\end{equation}
Then the maps and their inverses are equicontinuous.  Every
subsequential uniform limit is a \(K(c_*)\)-quasiconformal sphere
homeomorphism, with either orientation.
\end{lemma}

\begin{proof}
The two ring estimates below give equicontinuity and exclude collapsed
continua.  Annuli with shifted centers will then bound the linear
distortion of each limit.

\paragraph{Two ring estimates.}
Let \(E,F\) be the two
complementary continua of a sphere annulus.  If both have spherical
diameter at least \(u\), every essential closed curve in the annulus
has spherical length at least \(u/4\).  Otherwise its trace lies
in a disk of radius \(u/4\); each continuum has a point outside
that disk, so the curve is null-homotopic in the sphere minus those
two points, a contradiction.  Constant spherical density and
\eqref{conf:duality} give
\begin{equation}
 \EL_{\rm tr}(\Chat\setminus(E\cup F))
       \le 64\pi/u^2.
 \label{conf:ring-diameter-bound}
\end{equation}
If additionally \(d=\dist(E,F)\to0\) while both diameters stay
above \(u\), choose nearest points and center geodesic circles at
the point in \(E\).  For \(2d<t<u/4\), both continua meet the
circle of radius \(t\).  There is a circle arc joining them whose
interior is in the annulus.  Every admissible crossing density
therefore has integral at least one on that circle.  Cauchy--Schwarz
and polar integration imply
\begin{equation}
 \EL_{\rm tr}(\Chat\setminus(E\cup F))
       \le \frac{2\pi}{\log(u/(8d))}\longrightarrow0.
 \label{conf:ring-touching-bound}
\end{equation}
The density is extended by zero on the complementary continua when
performing this integration.

\paragraph{Equicontinuity and exclusion of collapsed continua.}
Fix \(\delta>0\) smaller than the pairwise limiting distances of
the target marks, and fix an image tolerance \(u>0\).  Choose
an integer \(k\) so large that
\(kc_*/2>64\pi/\min(u,\delta)^2\).  Around each source point
choose an outer coordinate disk omitting two limiting source marks
with clearance.  Inside it choose a stack of \(k\) rational
concentric annuli, with disjoint closures and a smaller central
disk.  Rational centers and radii can be chosen so that these
central disks cover the sphere; retain a finite subcover.  For this
finite collection, \eqref{conf:ring-hypothesis} gives a common
eventual index with all image moduli at least \(c_*/2\).  Series
addition and \eqref{conf:ring-diameter-bound} imply that every
retained central disk has image diameter less than \(u\): its
outer complementary continuum contains two target marks separated
by at least \(\delta\).  A Lebesgue number for the finite source
cover proves equicontinuity.  Only finitely many fixed annuli were
used for each \(u\).

Arzel\`a--Ascoli supplies a uniform limit \(f\), which is onto:
for each target point extract a convergent sequence of its
preimages.  Its fibers are connected.  Indeed, if a fiber split
into two disjoint nonempty compact sets, choose disjoint open
neighborhoods of these sets.  Their complement has image bounded
away from the fiber value.  For large \(n\), the connected
preimage under \(f_n\) of a sufficiently small disk about that
value would lie in the two neighborhoods and meet both, which is
impossible.

Suppose a fiber \(C\) is nontrivial.  Choose \(a\in\partial C\),
\(c\in C\setminus\{a\}\), and \(d\notin C\).  Take
\(b\notin C\) sufficiently close to \(a\).  A fixed rational
annulus can be chosen with \(a,b\) strictly inside its inner disk
and \(c,d\) strictly outside its outer disk.  The complementary
continua of its \(f_n\)-image have diameters bounded below by the
positive limiting distances from \(f(b)\) and \(f(d)\) to
\(f(a)=f(c)\), but their mutual distance tends to zero.  This
contradicts \eqref{conf:ring-touching-bound} and the lower bound
for this single fixed annulus.  Thus \(f\) is a homeomorphism.

\paragraph{Quasiconformality and inverse equicontinuity.}
We give the quantitative local criterion as well.  Work in finite
planar coordinates on a disk whose image avoids the target pole.
For a bounded planar annulus with through EL at least \(m>0\),
let \(e_0,e_1\) belong to its bounded complementary component and
\(w\) to the unbounded component.  The same circle integration
as above, with the unbounded continuum crossing every sufficiently
large circle, gives
\begin{equation}
 |w-e_i|\ge \theta(m)|e_1-e_0|\quad(i=0,1),
 \qquad \theta(m)=\tfrac12 e^{-2\pi/m}.
 \label{conf:relative-ring}
\end{equation}
For a fixed rational source annulus apply this with \(m=c_*/2\)
to \(f_n\), then pass to \(f\).  Approximation of centers,
radii, and points extends the inequalities to every sufficiently
small source annulus.  There is no common discrete threshold in
this approximation.

Put \(L(z,r)=\max_{|v-z|=r}|f(v)-f(z)|\) and define
\(l(z,r)\) by the corresponding minimum.  With
\(\theta=\theta(c_*/2)\), \eqref{conf:relative-ring} yields
\[
 l(z,2r)\ge\theta L(z,r).
\]
For \(|w-z|=2r\), let \(v=(z+w)/2\) and apply the same
inequality to the annulus centered at \(w\), where \(v\) lies
on its inner circle and \(z\) on its outer circle.  Provided
\(B(z,4r)\) is in the chart, this gives
\[
 |f(z)-f(v)|\ge\theta|f(v)-f(w)|,
 \qquad
 L(z,2r)\le(1+\theta^{-1})L(z,r).
\]
Consequently
\begin{equation}
 \frac{L(z,2r)}{l(z,2r)}
       \le \theta^{-1}(1+\theta^{-1}).
 \label{conf:linear-distortion}
\end{equation}
This is the metric criterion for planar quasiconformality
\cite[Theorem~1]{Cristea1989Quasiconformality}; see also
\cite[Theorem~2.5.3]{AstalaIwaniecMartin2009}. The
argument is precisely the lower-bound part of the proof of
\cite[Lemma~3.1, equations~(3.3)--(3.7)]{IvriiMarkovic2019}; the
shifted centers explain why one needs the all-center hypothesis.
One can also see the required Sobolev regularity directly:
\eqref{conf:linear-distortion} bounds \(L(z,r)^2\) by a constant
times \(\Area(f(B(z,r)))\).  Integrating squared coordinate
difference quotients, and using Fubini for the locally finite
measure \(V\mapsto\Area(f(V))\), gives bounded \(L^2\)
difference quotients on every compact subdomain.  Hence
\(f\in W^{1,2}_{\rm loc}\), and the metric criterion gives
bounded distortion.  Charts at the pole, or point removability,
finish the sphere assertion; a reflection handles the other
orientation.

Finally, failure of inverse equicontinuity would give target points
approaching one another whose source preimages stay separated.
Extract a uniform limit of the corresponding \(f_n\)'s and
convergent source preimages.  The limit would identify two distinct
points, contradicting the homeomorphism just proved.
\end{proof}

\subsection{The infinitesimal structure is round}
\label{conf:tensor-subsection}

The elimination of deterministic anisotropy by rotations has a parallel
in \cite[Section~4.4]{HoldenYuUniformization}, where determinism comes
from energy homogenization.  Here the conditional-product and field-germ
argument of Lemma~\ref{conf:germs} supplies that step.

\begin{lemma}[Conformal rigidity]
\label{conf:tensor}
Every homeomorphic limit in Lemma~\ref{conf:ring-criterion}, when
obtained from the local flag-surface arrays above, is conformal or
anticonformal.
\end{lemma}

\begin{proof}
Let \(f_n=\phi_n\circ H_n^{-1}\to f\) uniformly.  The image
of a local quadrilateral \(R\) has the same modulus as its actual
flag-surface preimage.  Strict inner and outer port versions
therefore determine the modulus of \(f(R)\): uniform convergence
of maps and inverses places the four sides between arbitrarily
small collars, with their marked sides in the correct order.
Conformal rectangle modulus is continuous under these Jordan
collar approximations.  The two transverse comparisons in
\eqref{conf:duality} supply the two sides of this determination.
Thus all these local moduli are readouts of the local EL arrays.
Global three-point normalization does not enter them, since
M\"obius postcomposition preserves modulus.

The distortion bound \eqref{conf:linear-distortion} implies that
\(f\) is differentiable almost everywhere
\cite[Lemma~2]{Cristea1989Quasiconformality}.  Its quasiconformality
also gives a nonzero Jacobian almost everywhere
\cite[Theorem~3.1.2]{AstalaIwaniecMartin2009}.
At a point \(z\) satisfying both conclusions put
\begin{equation}
 T(z)=\frac{Df(z)^{\mathsf T}Df(z)}{|\det Df(z)|}.
 \label{conf:tensor-definition}
\end{equation}
This positive symmetric tensor has determinant one and bounded
eccentricity.  It is a germ readout of the local moduli.  Here is
an explicit way to recover it.  For a unit vector \(v\) and
\(t>0\), let \(R_{v,t}\) have length one in direction \(v\)
and width \(t\).  Differentiability implies uniform convergence
of the rescaled images of \(z+rR_{v,t}\) to
\(Df(z)R_{v,t}\), for fixed \(t\).  If \(L\) is a nonsingular
linear map, constant density and straight-line flows give
\[
 \frac{(|Lv|-t|Lv^\perp|)_+^2}{t|\det L|}
 \le \EL_{\leftrightarrow}(LR_{v,t})
 \le \frac{|Lv|^2}{t|\det L|}.
\]
Multiplying by \(t\) and letting \(t\downarrow0\) recovers
\(v^{\mathsf T}T(z)v\).  Three fixed directions, with countable
sequences of \(r,t\) and strict covers, determine the tensor.
This proves both local readability and spatial measurability.

Lemma~\ref{conf:germs} makes \(T(z)\) deterministic for almost
every deterministic \(z\).  A coordinate rotation conjugates it
by the rotation matrix; the possible orientation sign of \(f\)
has disappeared from \eqref{conf:tensor-definition}.  Applying
the rotation part of that lemma to a quarter-turn gives
\(T=R^{\mathsf T}TR\).  For a symmetric two-by-two matrix this
sets the off-diagonal entries to zero and equates the diagonal
entries.  Since \(\det T=1\), we obtain \(T=I\) almost
everywhere.  The quasiconformal differential is therefore a
similarity almost everywhere, of one fixed orientation.  The
analytic characterization of quasiconformal maps
\cite[Theorem~2.5.4]{AstalaIwaniecMartin2009} makes \(f\)
conformal or anticonformal.  Isolated marked points are removable.
\end{proof}

\subsection{Actual marks, uniformization, and area}
\label{conf:normalization-subsection}

The remaining normalization uses the three fresh area samples of the
finite model.  We first record how these samples behave under the
reference homeomorphisms.

\begin{lemma}[Area marks under a converging coordinate]
\label{joint:marks}
Suppose $H_n:X_n\to\Chat$ are homeomorphisms and
$(H_n)_*m_n\Rightarrow\mu$ weakly, where $\mu$ is diffuse. Conditional
on the surface and exploration data used to choose $H_n$, let
$x_n^1,x_n^2,x_n^3$ be fresh independent samples from $m_n$. In
particular $H_n$ is chosen before these three marks. Then their images
converge jointly, together with the surface data, to three conditionally independent $\mu$ samples.
Almost surely these limiting points are distinct. The M\"obius map
sending the ordered triple to $0,1,\infty$ varies continuously with
the triple as long as the three points are distinct.
\end{lemma}

\begin{proof}
For continuous functions $f_1,f_2,f_3$ on $\Chat$, conditional
expectation of their product is
\[
 \EE\!\left[\prod_{j=1}^3 f_j(H_n(x_n^j))\,
        \middle|\,\text{surface data}\right]
   =\prod_{j=1}^3\int f_j\,\dd(H_n)_*m_n.
\]
Weak convergence passes the right side to
$\prod_j\int f_j\,\dd\mu$. Products of continuous test functions
determine the law on the compact product sphere, proving the joint
conditional assertion. Diffuseness gives distinctness. In a chart
avoiding the limiting points at infinity, the normalized M\"obius
map is the usual cross-ratio expression; its coefficients depend
continuously on the distinct triple. The resulting maps and their
inverses converge uniformly in spherical distance by compactness.
Other chart choices give the same conclusion when one point is
infinite.
\end{proof}

\begin{theorem}[Canonical uniformization comparison]
\label{conf:canonical}
Keep any joint extraction of the finite-map encoding and its local
observations.  Let \(x_n^1,x_n^2,x_n^3\) be the actual independent
\(m_n\)-samples used to define \(\phi_n\).  In the reference
coordinate these marks converge jointly to three conditionally
independent samples \(z^1,z^2,z^3\) of \(\mu_h\).  The maps
\(\phi_n\circ H_n^{-1}\) have uniformly convergent subsequences,
and every limit is the M\"obius or anti-M\"obius map taking
\((z^1,z^2,z^3)\) to \((0,1,\infty)\).  All retained local
observations use this same limiting coordinate and the same possible
orientation sign.
\end{theorem}

\begin{proof}
An original edge contributes four equal-mass flags and one
quadrangle.  Choosing a uniform quadrangle and then one of its
two matched tour steps gives a uniform step of the exploration.
By Proposition~\ref{top:projection} and Theorem~\ref{top:homeomorphism},
every point in any of its four
flags and either step representative have vanishing discrepancy
in the reference coordinate, uniformly over the map.  Hence, for
every continuous test function \(g\),
\[
 \int g\,\dd (H_n)_*m_n
       -\frac1{2n}\sum_{j=1}^{2n}g(\eta(j/(2n)))
       \longrightarrow0.
\]
The continuity of \(\eta\) and its quantum-area parameterization
identify the Riemann-sum limit as \(\int g\,\dd\mu_h\).
Repeated incidence labels cause no change: the four flags are
counted as occurrences.  Thus \((H_n)_*m_n\Rightarrow\mu_h\).

The maps \(H_n\) were chosen before the fresh samples.  Apply
Lemma~\ref{joint:marks} to this reference area convergence.  The
limiting marks have conditional law \(\mu_h^{\otimes3}\), jointly
with the retained map and local observations, and are almost surely
distinct.  Thus both the source marks and their prescribed target
images satisfy the normalization hypothesis of the ring criterion.

Conformal invariance identifies the through EL of
\((\phi_n\circ H_n^{-1})(A)\) with the actual surface EL for
the corresponding annulus.  Proposition~\ref{conf:macroscopic-el},
with strict collars, gives the hypothesis of
Lemma~\ref{conf:ring-criterion} simultaneously for the countable
annulus library.  That lemma gives uniform subsequential
convergence to a quasiconformal homeomorphism.  Lemma~\ref{conf:tensor}
makes it conformal or anticonformal, hence M\"obius or
anti-M\"obius on the sphere.  Uniform convergence and convergence
of the moving marks preserve their prescribed images.  These
three images specify the map once its orientation is fixed.
All extractions retained the same local arrays and topological
correspondence, so no observation acquires a different coordinate.
\end{proof}

\begin{corollary}[Canonical area and maximum mesh]
\label{conf:mesh-area}
In the joint coupling of Theorem~\ref{conf:canonical}, with
limiting normalization \(f\),
\[
 (\phi_n)_*m_n\Rightarrow f_*\mu_h,
 \qquad
 \max_T\diam_\rho\phi_n(T)\longrightarrow0.
\]
The measure on the right is the area of the quantum sphere in
the same three-marked coordinate; a common reflection is retained
when the correspondence reverses orientation.  In particular,
the maximum-mesh conclusion holds in probability along the full
sequence.
\end{corollary}

\begin{proof}
Write \(f_n=\phi_n\circ H_n^{-1}\).  Uniform convergence and
the reference area convergence prove the measure assertion.  If
\(\omega_f\) is a modulus of continuity of \(f\) on the sphere,
then
\[
 \max_T\diam_\rho\phi_n(T)
 \le2\|f_n-f\|_\infty+
       \omega_f\!\left(\max_T\diam_\rho H_n(T)\right)
       \longrightarrow0
\]
by \eqref{conf:topological-mesh}.  Every sequence has a further
joint extraction with these conclusions, which proves the stated
in-probability assertion about mesh.  The
three samples and \(\phi_n\) throughout are the prescribed
finite ones.
\end{proof}

\begin{corollary}[Area-typical root position]
\label{conf:root-area}
The exploration root has diffuse limiting area position in the
canonical coordinate.  This conditional sampling statement remains
valid jointly with any root-free map observations retained in the
same extraction.
\end{corollary}

\begin{proof}
Use the exact conditional rerooting statement in
Lemma~\ref{intro:root-flag}.  A uniform oriented-edge tail can be
coupled to a uniform flag by choosing one of its two sides, and
then to an independent uniform point of that flag.  The two
canonical positions differ by at most the maximum flag diameter.
Conditional on the root-free map data and the normalization
samples, this gives the same limiting continuous-test integrals
as sampling from \((\phi_n)_*m_n\).  Apply
Corollary~\ref{conf:mesh-area}.  The resulting area measure is
nonatomic.
\end{proof}
